\chapter{The Revolution in Perspective}

When Lenin in the April theses proclaimed that the revolution of February 1917 was not simply a bourgeois revolution, but marked a transition, led by workers and poor peasants, to the desired socialist revolution of the future, he made a sensitive response to the tumultuous conditions prevailing on his return to Petrograd. The Russian bourgeoisie, weak and backward in comparison with its western counterparts, possessed neither the economic strength nor the political maturity, neither the independence nor the inner coherence, necessary to wield power. On the other hand, the conception of an alliance between the proletariat and the bourgeoisie to complete the bourgeois revolution was pure myth. The proletariat, once it became an effective force, could not instal in power a bourgeois regime, whose function would be to exploit its labour. The bourgeoisie could not brook the alliance of a proletariat, whose eventual function would be to destroy it. When Lenin attempted to escape from this impasse by placing on the proletariat, supported by the poor peasants, the onus both of completing the bourgeois revolution and of starting the socialist revolution, he no doubt believed, not that he was abandoning the Marxist scheme of two separate and successive revolutions, but that he was adapting it to special conditions. But this solution, which became the programme of the October revolution, had its Achilles' heel. Marx had envisaged a socialist revolution developing on a foundation of capitalism and bourgeois democracy established by a previous bourgeois revolution.  In Russia this foundation was rudimentary or non-existent. Lenin looked forward to the building of socialism in an economically and politically backward country. The dilemma could be avoided only so long as it was assumed that the revolution was about to become international, that the European proletariat would also rise in revolt against its capitalist masters, and provide the conditions for an advance into socialism which Russia in isolation lacked. Socialism installed by revolution in a country where the proletariat itself was economically backward and numerically weak was not, and could not be, the socialism predicated by Marx and Lenin as the outcome of a revolution of the united proletariat of economically advanced countries. 
% 注：修正了 "bnce" 为 "once"，"la" 为 "a"，去除了 "[labour" 中的单侧括号。

From the first, therefore, the Russian revolution had a hybrid and ambiguous character. Marx noted that the embryo of bourgeois society had been formed within the matrix of the feudal order, and was already mature when bourgeois revolution installed it in the seats of power. It was assumed that something analogous would happen to socialist society before the victory of the socialist revolution. In one---but only one---respect this expectation was vindicated. Industrialization and technological modernization, which ranked high among the achievements of capitalist society, were also important pre-requisites of socialism. Long before 1914 the capitalist economies of the western world had begun to transcend the limits of small-scale production by individual entrepreneurs, and to substitute production by large-scale units which dominated the economic scene, and became involved willy-nilly in the exercise of political power. Capitalism itself was already blurring the line which separated economics from politics, paving the way for some form of centralized control and laying the foundations on which a socialist society could be built. 
% 注：去除了 "Ihybrid" 中的大写 "I"，去除了段末的乱码 "4" 和 "-f l|" 并与下文顺利拼接。

These processes reached a climax in the first world war. Study of the German war economy inspired Lenin's remark in the summer of 1917 that ``state monopoly capitalism is the fullest material preparation for socialism''; and a few weeks later he added, a little whimsically, that ``the material, economic half'' of socialism had been realized in Germany ``in the form of state monopoly capitalism''. The contradictions of capitalism had already produced, within the capitalist order, the embryo of the planned economy of the USSR. This fact has led some critics to describe what was achieved under Soviet planning as ``state capitalism''. Such a view seems untenable. Capitalism without entrepreneurs, without unemployment and without a free market, where no class appropriates the surplus value produced by the worker, and profits play a purely ancillary role, where prices and wages are not subject to the law of supply and demand, is no longer capitalism in any meaningful sense. Soviet planned economy was recognized everywhere as a challenge to capitalism. It was ``the material, economic half'' of socialism, and was a major outcome of the revolution. 
% 注：清除了 "Ttrrm trmrr rwr rrrrTr IWffT?" 这一长串由于页面污损造成的极端乱码。

If, however, it would be foolish to deny the title ``socialist'' to this achievement, it would be equally misguided to pretend that it constituted a realization of Marx's ``free association of producers'', or of the dictatorship of the proletariat, or of Lenin's transitional ``democratic dictatorship of workers and peasants''. Nor did it satisfy Marx's requirement that ``the emancipation of the workers must be the task of the workers themselves''. The Soviet industrial and agrarian revolutions plainly fell into a category of ``revolution from above'', imposed by the joint authority of party and state. The limitations of ``socialism in one country'' were plainly revealed. The vision of a trained and educated proletariat growing up within bourgeois society, like the bourgeoisie within feudal society, had not been realized---and, least of all, in backward Russia, where the working class was small, down-trodden, and unorganized, and had assimilated none even of the conditional freedoms of bourgeois democracy. The tiny core of class-conscious workers played a major part in the victory of the revolution. But the task of ordering and administering the broad territories incorporated in the Soviet republic called for a more complex and more sophisticated form of organization. The party of Stalin, a disciplined corps led by a small and devoted band of revolutionary intellectuals, stepped into the vacant place, and directed the policy of the regime by methods which, after Lenin's death, became more and more openly dictatorial, and less and less dependent on its proletarian base. Devices, first sparingly used amid the passions and atrocities of the civil war, were gradually elaborated into a vast system of purges and concentration camps. If the goals could be described as socialist, the means used to attain them were often the very negation of socialism. 
% 注：去除了 "^“the", "■ The" 等无意义的字符。修复了因断句错误产生的 "I, This does not mean"，将 "This does not mean..." 作为新段落的开头（见下段）。

This does not mean that no advance at all had been made towards the most exalted ideal of socialism---the liberation of the workers from the oppressions of the past, and the recognition of their equal role in a new kind of society. But progress was halting, and was broken by a series of set-backs and calamities, avoidable and unavoidable. After the ravages and shortages of the civil war, a brief respite ensued in which the standard of living of both workers and peasants slowly rose somewhat above the miserable levels of Tsarist Russia. During the decade which began in 1928, these standards once more contracted under the intense pressures of industrialization; and the peasant passed through the horrors of forced collectivization. Scarcely had recovery once more set in when the country was exposed to the cataclysm of a world war, in which the USSR was the target of Germany's most sustained and most devastating offensive on the continent of Europe. These terrifying experiences left their mark, material and moral, on Soviet life and on the minds of the Soviet leaders and the Soviet people. Not all the sufferings of the first half century of the revolution can be attributed to internal causes or to the iron hand of Stalin's dictatorship. 
% 注：去除了 "I ", "istil ^ ill" 这种由墨迹产生的断行乱码，并修正 "Alter" 为 "After"。

Yet, by the nineteen-fifties and nineteen-sixties, the fruits of industrialization, mechanization and long-term planning began to mature. Much remained primitive and backward by any western measurement. But standards of living substantially improved. Social services, including health services and primary, secondary and higher education, became more effective, and spread from the cities over most parts of the country. Stalin's most notorious instruments of oppression were dismantled. The pattern of life of ordinary people changed for the better. When the fiftieth anniversary of the revolution was celebrated in 1967, account could be taken of the magnitude of the advance. During the half century, the population of the USSR had increased from 145 millions to more than 250 millions; the proportion of town-dwellers had risen from less than 20 to more than 50 per cent. This was an immense increase in the urban population, most of the newcomers being children of peasants, and grandchildren or great-grandchildren of serfs. The Soviet worker, and even the Soviet peasant, of 1967 was a very different person from his father or grandfather in 1917. He could hardly fail to be conscious of what the revolution had done for him; and this out-weighed the absence of freedoms which he had never enjoyed or dreamed of. The harshness and cruelty of the regime were real. But so also were its achievements. 
% 注：修复了 "substan*- tially" 为 "substantially"。

Abroad, the immediate effect of the Russian revolution had been a sharp polarization of western attitudes between Right and Left. The revolution was a bugbear to conservatives, and a beacon of hope to radicals. Belief in this fundamental dichotomy inspired the foundation of Comintern. But, in the international revolution conceived by Marx and Lenin as a mass movement of the united European proletariat, no Marxist would have claimed a predominant role for the weak Russian contingent. When the European revolution failed to materialize, and when socialism in one country became the official ideology of the Russian party, the increasingly assertive demand to treat the USSR as the exemplar of socialist achievement, and Comintern as the repository of socialist orthodoxy, led to a new polarization between east and west within the Left. Communists and western social-democrats or socialists confronted one another, first as mistrustful allies, then as open enemies---a state of affairs misleadingly attributed in Moscow to the treachery of renegade leaders. It was a symptom of the rift that no common language could be found. International revolution as conceived in Moscow from 1924 onwards was a movement directed ``from above'' by an institution claiming to act in the name of the only proletariat which had made a victorious revolution in its own country; and the corollary of this re-orientation was the assumption, not only that the Russian leaders possessed a monopoly of knowledge and experience about the way in which a revolution could be made, but that the first and over-riding interest of international revolution was the defence of the one country where revolution had been effectively achieved. Both these assumptions, and the policies and procedures dictated by them, proved totally unacceptable to a majority of the workers of western Europe, who believed themselves far more advanced, economically, culturally and politically, than their backward Russian counterparts, and could not close their eyes to the negative aspects of Soviet society. Persistence in these policies merely brought discredit, in the eyes of western workers, on the authorities in Moscow, on the national communist parties subservient to them, and eventually on the revolution itself. 

Relations with the backward non-capitalist countries turned out quite differently. Lenin was the first to discover a link between the revolutionary movement for the liberation of the workers from capitalist domination in the advanced countries and the liberation of backward and subject nations from the rule of the imperialists. The identification of capitalism with imperialism was the fruitful theme of Soviet propaganda and policy almost everywhere in Asia, and enjoyed its most dramatic success in stimulating the Chinese national revolution in the middle nineteen-twenties. As the USSR consolidated its position, its prestige as the patron and leader of ``colonial'' peoples increased rapidly. It had achieved, through the process of revolution and industrialization, a spectacular accretion of economic independence and political power---an achievement worthy of envy and emulation. Outside Europe, even the exaggerated claims of Comintern made sense. The defence of the USSR, far from seeming an embarrassing excrescence on the programme of revolution, meant the defence of the most powerful ally of the backward countries in their struggle against the advanced imperialist countries. 
% 注：去除了 "lin the middle" 中的 l，修正了拼写 "excresence" 为 "excrescence"。

Nor did the methods which aroused revulsion in countries where the bourgeois revolution was a matter of history, and where strong workers' movements had grown up within the elastic framework of liberal democracy, prove seriously repugnant in countries where the bourgeois revolution was still on the agenda, where bourgeois democracy was an unsubstantial vision, and where no sizable proletariat yet existed. Where hungry and illiterate masses had not yet reached the stage of revolutionary consciousness, revolution from above was better than no revolution at all. While in the advanced capitalist world the ferment generated by the Russian revolution remained primarily destructive, and provided no constructive model for revolutionary action, in the backward non-capitalist countries it proved more pervasive and more productive. The prestige of a revolutionary regime which, largely through its own unaided efforts, had raised itself to the status of a major industrial Power, made it the natural leader of a revolt of the backward countries against the world-wide domination of western capitalism, which before 1914 had been virtually uncontested; and in this context the blots which tarnished its credentials in western eyes seemed irrelevant. Through the revolt of the backward non-capitalist world, the revolution presented a renewed challenge to the capitalist Powers, the potency of which is not yet exhausted. The Russian revolution of 1917 fell far short of the aims which it set for itself, and of the hopes which it generated. Its record was flawed and ambiguous. But it has been the source of more profound and more lasting repercussions throughout the world than any other historical event of modern times.